% =============================================================================
%  N1: Necessity — skeleton paper
%  STATUS: UNPROVEN. This file is a skeleton; the proof is not attempted here.
% =============================================================================
\documentclass[11pt]{article}

\usepackage[utf8]{inputenc}
\usepackage[T1]{fontenc}
\usepackage[expansion=false]{microtype}
\usepackage{amsmath, amssymb, amsthm, mathtools}
\usepackage[margin=1in]{geometry}
\usepackage{enumitem}
\usepackage[hidelinks]{hyperref}

\theoremstyle{plain}
\newtheorem{theorem}{Theorem}
\newtheorem{lemma}[theorem]{Lemma}
\newtheorem*{namedconjecture}{Conjecture}
\theoremstyle{definition}
\newtheorem{definition}[theorem]{Definition}

\newcommand{\X}{\mathbf{X}}

\title{Necessity of the Four-Slot Decomposition for Life-Producing Systems \\
       \large (UNPROVEN --- skeleton paper for N1)}
\author{Liam McGuire\thanks{Unaffiliated researcher, Seattle, USA. liammcg44@gmail.com}}
\date{\today}

\begin{document}
\maketitle

\begin{abstract}
This paper is a \textbf{skeleton} that records the statement of the
necessity conjecture N1 from the scaffolding
paper~\cite{emergent_systems_2026}, identifies the load-bearing
object (the decomposition functor $\mathrm{Decomp}$), and sketches the
operadic proof strategy. \textbf{No proof is attempted in this paper
as of its current revision.} The result is UNPROVEN.

N1 is the hardest of the framework's top-level conjectures. It
requires a constructive decomposition of arbitrary
$\mathcal{O}_W$-algebras into four-slot form, with proven essential
surjectivity on observer-observable behaviour.
\end{abstract}

\section{Statement (UNPROVEN)}

\begin{quote}
\emph{``N1 (necessity): for every life-producing pair
$(\mathcal{S}', O)$ in which $\mathcal{S}'$ is presentable as an
$\mathcal{O}_W$-algebra $A: \mathcal{O}_W \to \mathbf{Stoch}$,
there exists a 4-tuple system $\mathcal{S} = (\X, V, F, T)
\in \mathbf{Sys}$ and a slot-preserving morphism
$\mathcal{S} \to \mathcal{S}'$ that is essentially surjective on
$O$-observable behaviour.''}
\end{quote}
\noindent (Scaffolding paper~\cite{emergent_systems_2026},
\S\texttt{sec:n1}.)

\section{Plan of attack (no proof attempted)}

\begin{enumerate}[leftmargin=2em,itemsep=4pt]
  \item \textbf{Define $\mathrm{Decomp}$ (UNPROVEN).} Construct a
    decomposition functor
    $\mathrm{Decomp}: \mathrm{Alg}(\mathcal{O}_W) \to \mathbf{Sys}^{4\text{-slot}}$
    that takes an $\mathcal{O}_W$-algebra and returns a four-slot
    system. The construction must respect the operad's composition
    structure.
  \item \textbf{Generator check (UNPROVEN).} Apply Yau's finite-
    presentation theorem~\cite{yau2018operads} for $\mathcal{O}_W$
    (8 generators, 28 relations for the directed case; 6/17 for the
    undirected variant). Verify that $\mathrm{Decomp}$ respects each
    generator's action and each relation.
  \item \textbf{Essential surjectivity (UNPROVEN).} For every
    life-producing pair, $\mathrm{Decomp}$ should yield a 4-slot
    system whose $O$-observable behaviour matches. This is the
    behaviour-preservation condition.
  \item \textbf{Specialisation cases (UNPROVEN).} Instantiate
    $\mathrm{Decomp}$ for at least three canonical examples: a cell
    (chemical reaction network presentation), an immune-system
    trajectory (finite-state Markov chain presentation), and an
    ecological food web (coupled multi-substrate presentation).
\end{enumerate}

\section{Coordination}

N1 shares machinery with \texttt{op3\_topology\_universality/} (both
invoke the operad of wiring diagrams) and depends on
\texttt{c1\_closure\_unification/}
to make $\mathrm{Decomp}$ canonical on the viability side. The three
papers should be developed together.

\section{Status and follow-up}

\textbf{Current status: UNPROVEN.} \texttt{RESULTS.md} records N1 as
\texttt{open}. Likely the longest of the result papers; expect
substantial categorical machinery.

\bibliographystyle{plain}
\bibliography{references}

\end{document}
